% Part: turing-machines
% Chapter: undecidability
% Section: universal-tm

\documentclass[../../../include/open-logic-section]{subfiles}

\begin{document}

\olfileid{tur}{und}{uni}
\olsection{સર્વવ્યાપી ટ્યુરિંગ મશીનો}

\olref[enu]{sec}માં આપણે ચર્ચ્યું કે દરેક ટ્યુરિંગ મશીનનું
વર્ણન પૂર્ણાંકોની સાન્ત શ્રેણી વડે કરી શકાય છે. આ શ્રેણી
મશીનની અવસ્થાઓ, વર્ણમાળા, પ્રારંભિક અવસ્થા અને સૂચનાઓને
સંકેતિત કરે છે. આ બધાં વર્ણનોનો ગણ !!{enumerable} છે એ
પણ આપણે નોંધ્યું. આવાં વર્ણનોનો ગણ !!{denumerable} હોવાથી
~$\Nat$થી આ વર્ણનોમાં જતું કોઈ !!a{surjective} વિધેય છે.
દાખલા તરીકે, કેન્ટરની આડીઅવળી રીત વડે આવું
!!a{surjective} વિધેય મેળવી શકાય. તે બધાં ટ્યુરિંગ
મશીનોનાં (વર્ણનોનાં) પરિગણનાનો રસ્તો આપે છે. આવી કોઈ
એક પરિગણના નક્કી કરીએ, તો હવે $1$મું, $2$જું, \dots,
$e$મું ટ્યુરિંગ મશીન કહેવાનો અર્થ બને છે. આ અંકોને
\emph{સૂચકો} કહે છે.

\begin{defn}
આપણી નક્કી કરેલી પરિગણનામાં $M$ એ $e$મું ટ્યુરિંગ
મશીન હોય, તો $e$ને~$M$નો \emph{સૂચક} કહીએ છીએ.
$e$મા ટ્યુરિંગ મશીન માટે $M_e$ લખીએ છીએ.
\end{defn}

એક મશીનને એક કરતાં વધુ સૂચકો હોઈ શકે. દાખલા તરીકે,
~$M$નાં બે વર્ણનોમાં તેની સૂચનાઓની યાદીનો ક્રમ જુદો
હોઈ શકે; આવા જુદા વર્ણનોને જુદા સૂચકો મળશે.

અગત્યની વાત એ છે કે ટ્યુરિંગ મશીનોનાં વર્ણનોની પરિગણના
એવી રીતે આપી શકાય છે કે મશીનના સૂચકમાંથી~$M$નું વર્ણન
અસરકારક રીતે ગણી શકાય, અને મશીન~$M$ના વર્ણનમાંથી
તેનો કોઈ સૂચક અસરકારક રીતે ગણી શકાય. પછી ચર્ચ--ટ્યુરિંગ
માન્યતા મુજબ એવું ટ્યુરિંગ મશીન મેળવી શકાય છે જે
સૂચક~$e$વાળા ટ્યુરિંગ મશીનનું વર્ણન પાછું મેળવે અને
તે વર્ણનને આઉટપુટ તરીકે પોતાની ટેપ પર લખે. આ વર્ણન
~$\TMstroke$ના ખંડોની શ્રેણી હશે (જે~$M_e$નું વર્ણન
કરતી શ્રેણીના ધન પૂર્ણાંકો દર્શાવે છે).

આથી હવે સ્વાભાવિક રીતે પ્રશ્ન થાય: ટ્યુરિંગ મશીનોના
સૂચકો પરનાં કયાં વિધેયો પોતે ટ્યુરિંગ મશીન વડે સંગણનીય
છે? સૂચકોનાં કયાં ગુણધર્મોનો ટ્યુરિંગ મશીન વડે નિર્ણય
કરી શકાય? દાખલો લો: સૂચક~$e$ને સૂચક~$e$વાળા ટ્યુરિંગ
મશીનની અવસ્થાઓની સંખ્યા સાથે જોડતું વિધેય ટ્યુરિંગ
મશીન વડે સંગણનીય છે. આવું મશીન શું કરશે તે જુઓ:
તેની ટેપ પર $e$ જેટલા~$\TMstroke$નો એક ખંડ હોય ત્યારે
તે પહેલાં $e$માંથી મશીનનું વર્ણન કાઢશે. હવે ટેપ પર
વર્ણન~$\TMstroke$ના ખંડોની શ્રેણી વડે દર્શાવાયેલું છે.
આ શ્રેણીનો પહેલો !!{element} અવસ્થાઓની સંખ્યા છે.
તેથી હવે પહેલા~$\TMstroke$ના ખંડ સિવાયનું બધું ભૂંસીને
મશીન અટકી જાય એટલું જ કરવાનું રહે છે.

નીચેનું પરિણામ નોંધપાત્ર છે:

\begin{thm}\ollabel{thm:universal-tm} એવું \emph{સર્વવ્યાપી
  ટ્યુરિંગ મશીન}~$U$ છે કે ઇનપુટ $\tuple{e,n}$ પર શરૂ કરતાં
  \begin{enumerate}
    \item $M_e$ ઇનપુટ~$n$ પર અટકે ત્યારે અને માત્ર ત્યારે જ તે અટકે છે, અને
    \item $M_e$ આઉટપુટ $m$ સાથે અટકે, તો~$U$ પણ તેમ જ અટકે છે.
  \end{enumerate}
  આમ, $U$ એ $f\colon \Nat \times \Nat \pto \Nat$ વિધેય
  ગણે છે: $M_e$ ઇનપુટ~$n$ પર શરૂ કરીને આઉટપુટ~$m$
  સાથે અટકે ત્યારે $f(e,n) = m$, અને અન્યથા તે અવ્યાખ્યાયિત છે.
\end{thm}

\begin{proof}
  ખરેખર~$U$નું સંપૂર્ણ બાંધકામ આપવું વ્યવહારમાં લગભગ અશક્ય
  છે, કારણ કે તે અત્યંત જટિલ મશીન છે. પરંતુ તે કેવી રીતે
  કામ કરે છે તેની રૂપરેખા આપી, પછી ચર્ચ--ટ્યુરિંગ માન્યતાનો
  આધાર લઈ શકીએ. શરૂઆતમાં~$U$ની ટેપ પર $e$ જેટલા
  $\TMstroke$નો ખંડ, ત્યાર પછી $n$ જેટલા~$\TMstroke$નો
  ખંડ હોય છે. તે પહેલાં સૂચક~$e$માંથી મશીનનું વર્ણન
  કાઢીને ઇનપુટ~$n$ની જમણી બાજુ લખે છે. આમ મશીન~$M_e$ની
  સૂચનાઓનું વર્ણન કરતી સંખ્યાઓની યાદી મળે છે (એટલે
  ~$\TMblank$થી છૂટા પડેલા~$\TMstroke$ના ખંડો).
  પછી~$U$,~$M_e$ના વર્ણનની જમણી બાજુ,~$M_e$ની પ્રારંભિક
  અવસ્થાનો અંક અને સંખ્યા~$1$ ટેપ પર લખે છે. (ફરીથી,
  આ અંકો~$\TMstroke$ના ખંડો તરીકે એકકી સ્વરૂપમાં દર્શાવાય
  છે.) ત્યાર પછી તે ઇનપુટ ($n$ જેટલા~$\TMstroke$નો ખંડ)
  જમણી બાજુ નકલ કરે છે---પરંતુ દરેક $\TMstroke$ને
  ત્રણ $\TMstroke$ના ખંડથી બદલે છે (યાદ રાખો,
  $\TMstroke$ ચિહ્નનો અંક~$3$ છે; $\TMendtape$નો અંક
  $1$ અને $\TMblank$નો અંક $2$ છે). આ ખંડોની શ્રેણીના
  ડાબા છેડે (~$U$ની ટેપ પર $\TMblank$ ચિહ્નોથી છૂટા
  પડેલા ખંડોમાં) તે એક~$\TMstroke$ લખે છે, જે
  ~$\TMendtape$નો સંકેત છે.

  હવે~$U$ની ટેપ પર આ બધું છે: સૂચક~$e$, સંખ્યા~$n$,
  પ્રારંભિક અવસ્થાનો સંકેત-અંક (``વર્તમાન અવસ્થા''),
  પ્રારંભિક હેડ-સ્થાનનો અંક~$1$ (``વર્તમાન હેડ-સ્થાન''),
  અને ``ટેપ''ની પ્રારંભિક સામગ્રી (~$\TMblank$થી છૂટા
  પડેલા~$\TMstroke$ના ખંડોની શ્રેણી, જે~$M_e$નાં
  ચિહ્નોના સંકેત-અંકો---``ચિહ્નો''---દર્શાવે છે).

  હવે નીચેના પગલાં લઈને તે ઇનપુટ~$n$ પર શરૂ કરાયેલું
  $M_e$ શું કરે તેની અનુકૃતિ કરે છે:
  \begin{enumerate}
    \item ``વર્તમાન હેડ-સ્થાન''નો અંક~$k$ શોધો
    (શરૂઆતમાં તે~$1$ છે),
    \item ``ટેપ''ના $k$મા ખંડ પાસે જઈને ત્યાં કયું
    ``ચિહ્ન'' છે તે જુઓ,
    \item\ollabel{find-inst}%
    વર્તમાન ``અવસ્થા'' અને ``ચિહ્ન''ને મળતી સૂચના શોધો,
    \item ``ટેપ''ના $k$મા ખંડ પર પાછા જઈ ત્યાંનું
    ``ચિહ્ન''~$M_e$ જે ચિહ્ન લખે તેના સંકેત-અંકથી બદલો,
    \item વર્તમાન ``અવસ્થા''ની નોંધ જ્યાં છે ત્યાં હેડ
    લઈ જઈ એ અંકને નવી અવસ્થાના અંકથી બદલો,
    \item ``ટેપ-સ્થાન''ની નોંધ પાસે જઈ એક~$\TMstroke$
    ભૂંસો અથવા એક~$\TMstroke$ ઉમેરો (સૂચના અનુક્રમે
    ડાબે કે જમણે જવાનું કહે ત્યારે).
    \item ફરી કરો.\footnote{અહીં કેટલીક સૂક્ષ્મ મુશ્કેલીઓ
    છોડીને ચાલીએ છીએ. દાખલા તરીકે, ``વર્તમાન હેડ-સ્થાન''નો
    હિસાબ રાખતું ગણક વધારતી વખતે~$U$ને વધારે જગ્યા
    જોઈતી થઈ શકે---એ કિસ્સામાં આખી ``ટેપ'' જમણી બાજુ
    ખસેડવી પડશે.}
  \end{enumerate}
  ઇનપુટ~$n$ પર શરૂ કરાયેલું $M_e$ કદી ન અટકે, તો~$U$
  પણ કદી નથી અટકતું, એટલે તેનું આઉટપુટ અવ્યાખ્યાયિત છે.

  પગલું~\olref{find-inst}માં~$M_e$ના વર્ણનમાં વર્તમાન
  ``અવસ્થા''/``ચિહ્ન'' જોડ માટે કોઈ સૂચના નથી એવું જણાય,
  તો~$M_e$ અટકતું. આ થાય તો~$U$ પોતાની ટેપનો
  ``ટેપ''ની ડાબી બાજુનો ભાગ ભૂંસે છે. તે ત્રણ
  ~$\TMstroke$ના દરેક ખંડ માટે (જે~$M_e$ની ટેપ પરના
  એક~$\TMstroke$ને દર્શાવે છે) પોતાની ટેપના ડાબા છેડે
  એક $\TMstroke$ લખે છે અને ``ટેપ''ને ક્રમે ભૂંસે છે.
  આ પૂરું થાય ત્યારે~$U$ની ટેપ પર લંબાઈ~$m$વાળો
  ~$\TMstroke$નો એક જ ખંડ હોય છે.

  ``ટેપ'' પર ત્રણ~$\TMstroke$ના ખંડ સિવાયનું કંઈ
  $U$ને મળે, તો તે તરત અટકી જાય છે. આ કિસ્સામાં
  ~$U$ની ટેપ પર~$\TMstroke$નો એક જ ખંડ ન હોવાથી
  તેનું આઉટપુટ પ્રાકૃતિક સંખ્યા નથી, એટલે આ કિસ્સામાં
  $f(e,n)$ અવ્યાખ્યાયિત છે.
\end{proof}

\end{document}
